\zadanie{Elektronowy wkład do entropii tlenu}

Podstawowy poziom energii elektronowej cząsteczki \ce{O2} ma degenerację
$g_0 = 3$. Pierwszy wzbudzony poziom energii elektronowej ma degenerację
$g_1 = 2$ i leży $\tilde\nu = 7824$~cm$^{-1}$ powyżej poziomu podstawowego.
Wyższe poziomy elektronowe pomiń. Tlen traktuj jako gaz doskonały.

\begin{enumerate}[label=(\alph*)]
  \item Oblicz elektronowy wkład do entropii molowej tlenu, $S^\text{el}_m$,
    w temperaturze $T = 5000$~K.
  \item Oblicz $S^\text{el}_m$ w tej samej temperaturze, uwzględniając tylko
    poziom podstawowy. Porównaj oba wyniki.
\end{enumerate}

{\raggedright
Stałe: \mbox{$h = 6.626\times 10^{-34}$~J\,s},
\mbox{$c = 2.998\times 10^{10}$~cm/s},
\mbox{$k_B = 1.381\times 10^{-23}$~J/K},
\mbox{$R = 8.314$~J/(mol\,K)}.

Prawidłowa odpowiedź: (a)~$S^\text{el}_m = 10.9$~J/(mol\,K),
(b)~$S^\text{el}_m = 9.1$~J/(mol\,K).\par}

\begin{wskazowki}
\begin{itemize}
  \item Zapisz elektronową funkcję podziału jednej cząsteczki, $q^\text{el}$.
    Przyjmij energię poziomu podstawowego równą zeru i pamiętaj
    o degeneracjach.
  \item Cząsteczki nie oddziałują, więc dla $N$ cząsteczek
    $Z^\text{el} = (q^\text{el})^N$ (czynnik $1/N!$ poznany na wykładzie
    uwzględnia się w części translacyjnej).
  \item Entropia: $S^\text{el} = k_B\ln Z^\text{el} + E^\text{el}(T)/T$, gdzie
    $E^\text{el}(T) = k_BT^2\,\partial\ln Z^\text{el}/\partial T$.
  \item Energia wzbudzenia: $\varepsilon = hc\tilde\nu$.
\end{itemize}
\end{wskazowki}
